% BHU PYQ -- Auto-generated & error-corrected
\documentclass[11pt,a4paper]{article}

\usepackage[a4paper,top=2.5cm,bottom=2.5cm,left=2.8cm,right=2.8cm,headheight=14pt]{geometry}
\usepackage{amsmath,amssymb}
\usepackage[T1]{fontenc}
\usepackage[utf8]{inputenc}
\usepackage{lmodern}
\usepackage{microtype}
\usepackage{fancyhdr}
\usepackage{enumitem}
\usepackage{booktabs}
\usepackage{hyperref}
\usepackage{lastpage}
\usepackage{parskip}

\hypersetup{colorlinks=true,urlcolor=black,linkcolor=black,
  pdftitle={STA/STB-602: Demand Analysis, Income Distribution and Queuing Theory -- BHU PYQ},pdfauthor={Banaras Hindu University}}

\pagestyle{fancy}\fancyhf{}
\renewcommand{\headrulewidth}{0.4pt}\renewcommand{\footrulewidth}{0.4pt}
\fancyhead[L]{\small   \textit{Banaras Hindu University}}
\fancyhead[R]{\small   \textit{STA/STB-602: Demand Analysis \& Queuing Theory}}
\fancyfoot[C]{\small Page  \thepage\ of \pageref{LastPage}}


\newlist{parts}{enumerate}{2}
\setlist[parts,1]{label=(\alph*),leftmargin=*,itemsep=5pt,topsep=3pt}
\setlist[parts,2]{label=(\roman*),leftmargin=*,itemsep=3pt,topsep=2pt}

\newcommand{\pts}[1]{\hfill   \textit{[#1]}}


\begin{document}


\begin{center}
  {\large\bfseries BANARAS HINDU UNIVERSITY}\\[2pt]
  {\normalsize B.A./B.Sc. (Hons.) Semester VI Examination, 2022-23}\\[6pt]
  \rule{0.8\linewidth}{0.5pt}\\[6pt]
  {\large\bfseries Paper No. STA/STB-602: Demand Analysis, Analysis of Income Distribution and Queuing Theory}\\[4pt]
  {\normalsize Time: 3 Hours\hfill Maximum Marks: 70}
\end{center}
\rule{\linewidth}{0.4pt}
\smallskip



\noindent   \textit{Instructions: Answer any   \textbf{five} questions. All questions carry equal marks. Symbols have their usual meanings.}
\bigskip

\medskip


\noindent   \textbf{Question 1.}

\begin{parts}
  \item Write a note on the equilibrium price. Find the equilibrium price and quantity for crude oil if the demand and supply equations are $Q_d = 500 - 4p$ and $Q_s = 100 + 2p$, respectively. Also sketch a free-hand curve and find the equilibrium price and quantity graphically. \pts{7}
  \item Define a queuing system and its important components. \pts{7}
\end{parts}

\medskip\rule{\linewidth}{0.2pt}

\medskip


\noindent   \textbf{Question 2.}

\begin{parts}
  \item Describe Pigou's method (with assumptions and limitations) for estimating elasticity from both time series data and family budget data. \pts{7}
  \item Consider a competitive market for which the quantities demanded and supplied (per year) at various prices are given as follows:
  
\begin{center}
  
\begin{tabular}{c|cccc}
   oprule
  Price (Rs.) & 600 & 800 & 1000 & 1200 \\
  \midrule
  Demand (million) & 22 & 20 & 18 & 16 \\
  Supply (million) & 14 & 16 & 18 & 20 \\
  ottomrule
  \end{tabular}
  \end{center}
  
\begin{parts}
    \item Calculate the price elasticity of demand when price is Rs.\ 800.
    \item Calculate the price elasticity of supply when price is Rs.\ 800.
    \item What are the equilibrium price and quantity?
    \item Suppose the government sets a price ceiling of Rs.\ 800. Will there be a shortage, and if so, how large will it be?
  \end{parts}\pts{7}
\end{parts}

\medskip\rule{\linewidth}{0.2pt}

\medskip


\noindent   \textbf{Question 3.}

\begin{parts}
  \item For an exponential income distribution $dF = c\,e^{-cx}\,dx$, $c > 0$, $0 < x < \infty$, show that the Lorenz coefficient of concentration is equal to $\frac{1}{2}$. \pts{7}
  \item The data below relates to the distribution of urban households in India. Fit a Pareto distribution to the data and, with the estimated parameters, derive the Lorenz concentration coefficient and the mean income:

  
\begin{center}
  \small
\begin{tabular}{l|c}
   oprule
  Income Class (Rs.) & Households (\%) \\
  \midrule
  Under 500 & 13.6 \\
  500 -- 999 & 28.9 \\
  1000 -- 1999 & 32.5 \\
  2000 -- 2999 & 10.6 \\
  3000 -- 3999 & 5.6 \\
  4000 -- 4999 & 3.1 \\
  5000 -- 5999 & 1.7 \\
  6000 -- 7999 & 1.7 \\
  8000 -- 9999 & 0.7 \\
  10000 -- 14999 & 0.8 \\
  15000 -- 24999 & 0.5 \\
  25000 and above & 0.3 \\
  ottomrule
  \end{tabular}
  \end{center}\pts{7}
\end{parts}

\medskip\rule{\linewidth}{0.2pt}

\medskip


\noindent   \textbf{Question 4.}

\begin{parts}
  \item Prove that in a Poisson arrival process, the inter-arrival time distribution follows the exponential distribution. \pts{7}
  \item Define a machine-servicing model in queuing theory. Derive the expression for the probability of $n$ customers in the system, $P_n$, for the finite-capacity $(M/M/1): (N/ \text{FCFS})$ model. \pts{7}
\end{parts}

\medskip\rule{\linewidth}{0.2pt}

\medskip


\noindent   \textbf{Question 5.}

\begin{parts}
  \item Define the cumulative distribution of the waiting time for a customer who has to wait in the $(M/M/1): (\infty/ \text{FCFS})$ queuing model. \pts{7}
  \item A one-person barber shop can accommodate a maximum of 5 people at a time (4 waiting, 1 getting a haircut). Customers arrive according to a Poisson distribution with mean rate 5 per hour. The barber cuts hair at an exponential rate of 4 per hour. Find:
  
\begin{parts}
    \item What percentage of time is the barber idle?
    \item What fraction of potential customers is turned away?
    \item What is the expected number of customers waiting for a haircut?
    \item How much time can a customer expect to spend in the barber shop?
  \end{parts}\pts{7}
\end{parts}

\medskip\rule{\linewidth}{0.2pt}

\medskip


\noindent   \textbf{Question 6.}

\begin{parts}
  \item For the case of $s$ channels with Poisson arrivals at rate $\lambda$ and exponential service at rate $\mu$, show that the probability that an arrival has to wait is:
  \[
    P( \text{wait}) = \frac{(\lambda/\mu)^s}{s!\left(1 - \frac{\lambda}{s\mu}\right)} \cdot P_0
  \]\pts{7}
  \item At a health centre, one patient arrives every $\frac{1}{2}$ hour. The average treatment time for a patient is 20 minutes. If the hospital has 2 doctors, find:
  
\begin{parts}
    \item Probability that both doctors are idle
    \item Probability that exactly one doctor is idle
    \item Expected number of patients waiting for a doctor
    \item Expected waiting time of a patient
  \end{parts}\pts{7}
\end{parts}

\medskip\rule{\linewidth}{0.2pt}

\medskip


\noindent   \textbf{Question 7.} Write short notes on any   \textbf{three} of the following:\pts{14}

\begin{parts}
  \item Curves of concentration (Lorenz curve)
  \item Pareto's law of income distribution
  \item Giffen's paradox
  \item Transient and steady states in queuing theory
\end{parts}

\vfill
\rule{\linewidth}{0.4pt}

\begin{center}\small
  \href{https://bhu-exam.vercel.app/}{Click here to visit our website}\\[4pt]
  \href{https://me-aryan.vercel.app/}{Click here to visit our contributor}\\[4pt]
  \href{https://quashan-me.vercel.app/}{Click here to visit our contributor}
\end{center}

\end{document}